%! Author = lili
%! Date = 5/5/26

\defin{Prolog}

\section{Why Prolog for KR?}
Prolog is usually taught as a programming language, but it is also a knowledge representation language.

We can view programs (also prolog) from two perspectives which both apply.

\paragraph{Declarative / KR view} is that a program is a set of logical definitions. You describe
what is true, not how to compute it.

\paragraph{Procedural view} is that a program is a set of instructions that Prolog executes to
answer queries step by step.

Since both views apply to the same program, this means the same program is both
knowledge representation and executable reasoning.
% todo häh und was sagt uns das genau?

\komm{facts represent anything that could be true}

When you write a Prolog program as a set of logical definitions, the same rules can be
used for many tasks, not just answering a fixed query.

% todo frage sind das die Rules odder types of rules?

% TODO entwerder defis oder als explanation wie es finktioniert
\defin{deduction}
\defin{abduction}
\defin{containment testing}
\defin{view differentiation}
\defin{query folding}

None of these are straightforward if you treat the program as pure code.
% todo why? dont quite understand it

\section{Prolog and First-Order Logic}
Basically Prolog works the same as fol, we have predicates and rules, but it is not the same!
Prolog was designed as a direct computational realization of first-order predicate logic. Each concept in FOL has an exact counterpart in Prolog:
\begin{table}[H]
    \centering
    \begin{tabular}{ll}
        \textbf{FOL concept} & \textbf{Prolog equivalent} \\
        \hline
        Predicate symbol & predicate name, e.g.\ parent, mortal \\
        $\forall x$ & variable (any uppercase name) \\
        Conjunction $\land$ & comma `,` in rule body \\
        Implication $A \rightarrow B$ & rule \texttt{B :- A.} \\
        Ground atom (no variables) & fact \\
        Horn clause & any Prolog clause \\
    \end{tabular}
\end{table}
This is why reading a Prolog program declaratively is reading logic.
% todo was bedeutet "reading a Prolog program declaratively"??? wie kann man es anders lesen?

\komm{F8: ":-" ist a headsign(neck)}

\komm{"? - mortal ( socrates ) ." ist the query}

\defin{horn-clause}

% todo verallgemeinerung F 8  wie man das übersetzt / formalisiert

\begin{exa}
    F8 % todo farblich codiert gegenüber stellen
\end{exa}

FOL is descriptive but not executable and therefore needs external theorem provers. Prolog provides built-in
inference and supports automatic query answering.
% todo diesen satz nochmal genau auseonander nehmen

% todo was genau bedeutet escriptive but not executable? wie muss man sich das in der Realität vorstellen?

\section{Prolog Syntax}
Before writing programs, we need to know what Prolog data looks like. Everything in
Prolog is a term. There are four kinds:
\defin{atom}

\defin{compound term}

\begin{exa}
    clauses F 12 % todo
\end{exa}

\defin{prolog-number}
\defin{prolog-variable}

% todo frage four kinds


A Prolog program consists of clauses. There are three kinds:
% todo frage typen von Klauseln
\defin{prolog-fact}

\begin{exa}
    clauses F 11 % todo
\end{exa}

\defin{prolog-rule}

\begin{exa}
    clauses F 11 % todo
\end{exa}

\defin{prolog-query}

\begin{exa}
    clauses F 11 % todo
\end{exa}

\begin{notation}
    Comma , in a rule body means AND. The neck :- means IF
\end{notation}

% todo fragen zu wie schreibt man das
\defin{Identity of a functor} % F 12

\section{A Running Example: Family Relations}
To learn Prolog syntax and reasoning, we will use one concrete example throughout the
next several slides: a family tree.

The family tree contains these people: % todo

% todo schritte F 15 + erst generalisiert aufschreiben, dann



\section{How Prolog Executes: Matching and Search}

\section{Declarative vs. Procedural Meaning}

\section{Lists}

% zu f30 todo bsp von Tafel

We need to specify head and tail for every list if we want to do smth like append ([ H | T1 ] , List2 , [ H | T3 ]) : -

\begin{exa} Extra
\[
    A =  [ \underbrace{1,2}_{H}, \underbrace{3,4}_{T} ] \\
    B =  [ \underbrace{d}_{H}, \underbrace{e,f}_{T} ]
\]
If we now do that:
\[
    append([\underbrace{H_1}_{1} \mid  \underbrace{T_1}_{3}], B, [\underbrace{H_2}_{2} \mid  \underbrace{T_2}_{4}],) :-
\]
we get
\[
    [1,3,d,e,f,2,4]
\]
\end{exa}

\komm{built-in =  pre defined}

\section{Recursion and Arithmetic}



\section{Clause Ordering and Loops}

\section{The Monkey and Banana Problem}

\section{Backtracking, Cuts, and Negation}

\section{Horn Clauses and Proof}

\section{Operators}

% todo das mit den predecessors ist whack beschrieben, das muss ich mir genauer anschauen und nachfragen (auf Foto)

%  F 12 - 30 nochmal durchgehen wegen defis usw.